\documentclass{article}

\usepackage{report_style}
\usepackage[utf8]{inputenc}
\usepackage[T1]{fontenc}
\usepackage{amsmath}
\usepackage{booktabs}
\usepackage{float}
\usepackage{graphicx}
\usepackage{microtype}
\usepackage{tabularx}
\usepackage{xcolor}
\usepackage{hyperref}
\usepackage{url}

\newcolumntype{Y}{>{\raggedright\arraybackslash}X}
\newcommand{\code}[1]{\texttt{#1}}
\hypersetup{
  colorlinks=true,
  linkcolor=black,
  citecolor=blue!55!black,
  urlcolor=blue!55!black,
  pdftitle={Nano-ESMC: A Minimal, Hill-Climbable ESMC Training Reproduction},
  pdfauthor={Lumin Science}
}

\title{Nano-ESMC:\\
A Minimal, Hill-Climbable ESMC Training Reproduction}

\author{%
  Lumin Science\\
  \texttt{https://luminscience.io}
}

\begin{document}
\raggedbottom
\maketitle

\begin{abstract}
Nano-ESMC is a compact, end-to-end reconstruction of ESMC-style protein
language-model training. It connects a public 665,970,495-protein,
evaluation-decontaminated corpus to a fixed four-GPU training contract and
frozen biological evaluation. Full long-range contact precision at $L$
(P@L) on 20,775 RCSB Protein Data Bank chains is the primary model-selection
axis; held-out masked-language-model loss and the four-task P-CORE
representation panel are supporting measurements. Under the current one-hour
four-L40S contract, the original ESMC-300M-class baseline reaches 0.0925 P@L
and the first retained AutoResearch candidate reaches 0.0960. Released
ESMC-300M, 600M, and 6B checkpoints establish the external scale ladder for
P-CORE, while a source-qualified P-CORE Next alpha campaign records which new
tasks discriminate scale and which fail qualification. This report presents
the data pipeline once, the evaluation contract once, and all baseline
evidence in one section.
\end{abstract}

\section{Scope and fixed contract}

The goal is to reproduce the essential ESMC sequence-model training path in a
form that is small enough to inspect and strict enough to hill climb
\citep{candido2026language}. A candidate may change model internals,
optimization, batching, kernels, or schedules, but it cannot change the data,
compute budget, evaluation examples, or metric after observing a result.

\begin{table}[H]
  \centering
  \small
  \begin{tabularx}{\textwidth}{@{}p{0.18\textwidth}Y@{}}
    \toprule
    Object & Current contract \\
    \midrule
    Model class & Sequence-only ESMC-300M encoder trained from scratch; no structure-conditioned input. \\
    Data & Immutable public UniRef90, MGnify, and OMG/IMG release with exact and homology decontamination against the complete protected evaluation union. \\
    Compute & Four NVIDIA L40S GPUs and 3,600 seconds of synchronized training-loop time. \\
    Primary selection & Full long-range contact P@L on the frozen 20,775-chain population. \\
    Supporting axes & Held-out sequence-mean MLM loss and the four-task P-CORE frozen-representation panel. \\
    Promotion & Strict P@L improvement at matched compute, no unexplained trusted-task regression, then clean rebuilds and repeated seeds before a release claim. \\
    \bottomrule
  \end{tabularx}
  \caption{The public Nano-ESMC experiment contract.}
  \label{tab:contract}
\end{table}

The current repository supports Stage-1 training at context length 512.
Longer-context Stage 2, larger local models, and week-scale compute projections
are outside the present claim. Historical experiments remain evidence about
scaling, not alternate active contracts.

\section{Data: one corpus and one processing pipeline}
\label{sec:data}

\subsection{Sources and the main gap to ESMC}

ESMC reports training pools from UniRef, MGnify, and a July 2023 Joint Genome
Institute (JGI) snapshot \citep{candido2026language,suzek2015uniref,
richardson2023mgnify}. Nano-ESMC reconstructs the same source roles with
UniRef90 2023\_02, MGnify 2023\_02, and the IMG/JGI-labelled portion of the
public OMG corpus \citep{cornman2024omg}. OMG/IMG is a surrogate for the JGI
role, not the authors' exact snapshot.

\begin{table}[H]
  \centering
  \scriptsize
  \resizebox{\textwidth}{!}{%
  \begin{tabular}{@{}lrrr@{}}
    \toprule
    Source role & ESMC 70\%-identity clusters & Nano-ESMC representatives & Nano / ESMC \\
    \midrule
    UniRef & approximately 83M & 92,230,941 & approximately 111\% \\
    MGnify & approximately 372M & 348,135,082 & approximately 94\% \\
    JGI role / OMG-IMG surrogate & approximately 2B & 324,923,979 & approximately 16\% \\
    \midrule
    Total & approximately 2.455B & 765,290,002 & approximately 31\% \\
    \bottomrule
  \end{tabular}}
  \caption{Stage-aligned comparison with the cluster counts reported by ESMC.
  The roughly 1.68B-representative JGI-role gap dominates the difference.}
  \label{tab:source-gap}
\end{table}

\subsection{Processing and release accounting}

The complete path is fixed in seven steps:

\begin{enumerate}
  \item pin every upstream object and its digest;
  \item normalize sequences and assign OMG accessions to one source role;
  \item reject sequences shorter than 60 amino acids or with more than 20\%
        ambiguous residues;
  \item collapse exact sequences while retaining source membership;
  \item cluster each source view independently at 70\% identity and 80\%
        shorter-sequence coverage with MMseqs2 Linclust
        \citep{steinegger2018linclust};
  \item remove exact matches and homologs to the protected evaluation union;
        and
  \item assign cross-source ownership, select validation rows, write
        deterministic Parquet shards, and independently verify every digest.
\end{enumerate}

\begin{table}[H]
  \centering
  \scriptsize
  \setlength{\tabcolsep}{3.0pt}
  \resizebox{\textwidth}{!}{%
  \begin{tabular}{@{}lrrrrrr@{}}
    \toprule
    Source & Original records & Quality eligible & Distinct in source & 70\% representatives & After eval. decontam. & Final training \\
    \midrule
    UniRef90 & 170,669,877 & 165,884,293 & 165,884,293 & 92,230,941 & 74,181,938 & 74,175,974 \\
    MGnify & 729,215,663 & 611,817,034 & 611,788,129 & 348,135,082 & 330,444,705 & 328,949,335 \\
    OMG/IMG & 3,280,269,924 & 2,085,746,535 & 963,673,186 & 324,923,979 & 282,202,559 & 262,845,186 \\
    \midrule
    Total source views & 4,180,155,464 & 2,863,447,862 & 1,741,345,608 & 765,290,002 & 686,829,202 & 665,970,495 \\
    \bottomrule
  \end{tabular}}
  \caption{The single corpus funnel. Source-view distinct counts can overlap
  because one exact sequence may carry more than one source membership.
  Cross-source ownership is resolved before the final training count.}
  \label{tab:data-funnel}
\end{table}

The protected union contains 317,000 unique evaluation sequences, including
all probe-fit, validation, gallery, test, blocked-candidate, and contact rows
known at release-build time. It includes the 20 contact-probe chains and all
20,775 contact-evaluation chains from the frozen 2024-02-28 RCSB PDB snapshot.
The pipeline removes 21,653 exact matches and 78,439,147 additional homologs
at least 30\% identity, at least 80\% bidirectional coverage, and
$E\text{-value}\leq10^{-3}$ using MMseqs2 \citep{steinegger2017mmseqs2}.
After decontamination, 20,844,431 non-owner source memberships, 1,988
out-of-range records, and 12,288 validation rows are excluded to reach the
final training release.

The immutable public artifact contains 565 training shards,
665,970,495 proteins, and 151,304,238,405 residues. Each source contributes
one additional 4,096-protein validation shard. The complete train/validation
release occupies 109,661,312,410 compressed bytes. These values define the
supported corpus. The earlier 9M-protein materialization appears only in the
historical baseline table in Section~\ref{sec:evidence}.

\section{Evaluation}
\label{sec:evaluation}

\subsection{Three reported axes}

\begin{table}[H]
  \centering
  \small
  \begin{tabularx}{\textwidth}{@{}p{0.20\textwidth}p{0.18\textwidth}Y@{}}
    \toprule
    Axis & Role & Definition \\
    \midrule
    Held-out MLM & Diagnostic & Sequence-mean negative log-likelihood and perplexity on 12,288 hash-disjoint validation representatives. \\
    P-CORE & Representation panel & Four frozen-encoder tasks spanning evolution, local structure, localization, and fitness. \\
    Long-range contact P@L & Primary selection & Mean top-$L$ precision over the frozen 20,775-chain population using all-layer/all-head attention features. \\
    \bottomrule
  \end{tabularx}
  \caption{Evaluation axes and their decision roles. A lower MLM loss does not
  override a contact P@L regression.}
  \label{tab:evaluation-axes}
\end{table}

\subsection{Contact P@L}

All contact structures and labels come from the 2024-02-28 RCSB Protein Data
Bank snapshot \citep{berman2000pdb}. Following the ESM contact protocol and
the ESMC long-range definition \citep{rao2021unsupervised,
candido2026language}, all attention heads and layers are symmetrized and fed
to a logistic probe. Sixteen structures fit candidate probes, four select
regularization, the selected probe is refit on all 20, and 20,775 chains
supply the final score. A true contact has
C$\beta$ distance below 8\AA{} (C$\alpha$ for glycine), long range means
sequence separation at least 24, and the per-chain metric is precision among
the top $L$ predicted pairs.

The snapshot yields 216,478 source structures and 781,077 extracted protein
chains before filtering. Chains are reduced at 40\% sequence identity and 80\%
shorter-sequence coverage, truncated to 510 residues, and required to contain
at least $L$ eligible true contacts. The final population contains 20,775
chains representing 20,758 unique sequences. ESMC reports the same snapshot
date, population size, and construction rules, but does not publish its
ordered chain manifest or raw-file digests. The two evaluations are
protocol-matched, not proven item-for-item identical.

\subsection{P-CORE}

The encoder remains frozen. Protein tasks mean-pool final-layer residue
vectors; secondary structure uses residue vectors directly. A standardizer
and a fixed low-capacity logistic or ridge probe are fit on probe-training
data, regularization is selected on validation, and the selected probe is
applied once to the untouched test set. Bootstrap intervals resample frozen
test predictions by biological unit and do not refit the encoder or probe.
The task sources are DeepSF/TAPE for remote homology
\citep{hou2018deepsf,rao2019tape}, NetSurfP-2.0/CB513 for secondary structure
\citep{klausen2019netsurfp,cuff1999evaluation}, DeepLoc2 for localization
\citep{thumuluri2022deeploc2}, and FLIP2 for fitness
\citep{didi2026flip2}.

\begin{table}[H]
  \centering
  \small
  \setlength{\tabcolsep}{3pt}
  \begin{tabularx}{\textwidth}{@{}p{0.18\textwidth}p{0.25\textwidth}p{0.20\textwidth}YY@{}}
    \toprule
    Task & Fit / validation / test & Frozen readout & Primary metric & Unit \\
    \midrule
    Remote homology & 12,312 / 736 / 718 proteins & mean + multinomial logistic & balanced accuracy & family group \\
    Secondary structure & 8,678 / 2,170 / 513 proteins & residue + logistic & residue macro-F1 & protein \\
    DeepLoc2 & 28,303 proteins in five rotating folds & mean + multilabel logistic & macro average precision & fold / protein \\
    FLIP2 Hydrophobic Core & 9,974 / 2,493 / 12,468 variants & mean + ridge & Spearman correlation & variant group \\
    \bottomrule
  \end{tabularx}
  \caption{The four trusted P-CORE tasks. Each raw metric is null-normalized;
  the score is 100 times their geometric mean. A task at or below its null
  invalidates the aggregate.}
  \label{tab:pcore-contract}
\end{table}

Enzyme Commission and the reconstructed Human PPI benchmark remain visible
diagnostics, but neither can affect P-CORE or model selection. Their released
model behavior is reported once in Section~\ref{sec:evidence}.

\section{Baseline and qualification evidence}
\label{sec:evidence}

\subsection{Current one-hour AutoResearch baseline}

\begin{table}[H]
  \centering
  \scriptsize
  \setlength{\tabcolsep}{3.4pt}
  \resizebox{\textwidth}{!}{%
  \begin{tabular}{@{}lrrrrrrrrr@{}}
    \toprule
    Model & P@L & Delta & Train loss & Valid. loss & Steps & Tokens (M) & Params. (M) & Peak GB & Train h \\
    \midrule
    Original ESMC & 0.0925 & -- & 2.723 & 2.699 & 5,883 & 358 & 333 & 35 & 1 \\
    AutoResearch-Codex-Round1 & \textbf{0.0960} & \textbf{+0.0036} & \textbf{2.720} & 2.704 & 5,682 & 346 & 333 & 37 & 1 \\
    \bottomrule
  \end{tabular}}
  \caption{Compute-matched Stage-1 runs on four L40S GPUs. P@L is the keep or
  reject metric; losses are required diagnostics.}
  \label{tab:autoresearch}
\end{table}

The retained candidate adds learned residual/input routing, parameter-free
transformer RMSNorm, depth-scaled residual initialization, and a linear
learning-rate cooldown over the final 20\% of training. Its improvement is a
single fixed-checkpoint result, not a multi-seed release claim.

\subsection{Released ESMC checkpoints and quarantined diagnostics}

\begin{table}[H]
  \centering
  \scriptsize
  \setlength{\tabcolsep}{4pt}
  \begin{tabularx}{\textwidth}{@{}Yrrrl@{}}
    \toprule
    Metric & ESMC-300M & ESMC-600M & ESMC-6B & Role \\
    \midrule
    Remote homology balanced accuracy & 0.1159 & 0.1152 & 0.1186 & P-CORE \\
    Secondary-structure macro-F1 & 0.8315 & 0.8407 & 0.8780 & P-CORE \\
    DeepLoc2 macro average precision & 0.6442 & 0.6568 & 0.6942 & P-CORE \\
    FLIP2 Hydro Spearman & 0.4132 & 0.4276 & 0.4616 & P-CORE \\
    \textbf{P-CORE} & \textbf{37.4847} & \textbf{38.1553} & \textbf{40.6011} & reference panel \\
    \midrule
    Enzyme Commission macro average precision & 0.7174 & 0.7096 & 0.0237 & quarantined \\
    Human PPI average precision & 0.8155 & 0.8026 & 0.8345 & quarantined \\
    \midrule
    ESMC paper contact P@L-LR & $0.552\pm0.002$ & $0.589\pm0.002$ & $0.725\pm0.002$ & published \\
    Nano-ESMC full 20,775-chain P@L & 0.5387 & 0.5803 & pending & local protocol \\
    \bottomrule
  \end{tabularx}
  \caption{One consolidated released-model reference table. P-CORE values are
  local recomputations with the frozen probes. Contact paper values come from
  ESMC; local contact values use the Nano-ESMC chain manifest.}
  \label{tab:released-baselines}
\end{table}

P-CORE is ordered across the three released scales, although remote homology
is nearly flat. EC collapses at 6B despite clean split and embedding-integrity
checks. Human PPI has only 237 test pairs, is non-monotone across scale, and
was nearly matched by a severely undertrained local model. These behaviors
justify quarantine rather than post hoc repair. The full 6B local contact
evaluation remains pending and is not replaced by the completed 1,024-chain
diagnostic.

\subsection{Historical training-duration check}

\begin{table}[H]
  \centering
  \small
  \begin{tabular}{@{}lrrrrl@{}}
    \toprule
    Historical run & Steps & Model tokens & P-CORE & Contact P@L & Decision \\
    \midrule
    Local ESMC-300M, 4 h & 20,971 & 1.290B & 25.4609 & 0.1013 & baseline \\
    Local ESMC-300M, 16 h & 84,000 & 5.166B & 30.6609 & 0.1636 & improved, not promoted \\
    Released ESMC-300M & -- & -- & 37.4847 & 0.5387 & external reference \\
    \bottomrule
  \end{tabular}
  \caption{Historical four-A100 duration experiment. Both local runs used the
  superseded 9M-protein materialization, so they are retained as scaling
  evidence and are not baselines for the current public corpus.}
  \label{tab:historical-duration}
\end{table}

All four trusted raw tasks improved from four to sixteen hours. Contact P@L
increased by 0.0623, with a fixed-checkpoint paired 95\% interval of
[0.0614, 0.0632]. The runs used different seeds, so that interval does not
isolate training duration from initialization. The sixteen-hour checkpoint
remained below released ESMC-300M on P-CORE and contact and was not promoted.

\subsection{P-CORE Next alpha qualification}

The expansion campaign evaluated identical manifests with released ESMC-300M,
600M, and 6B checkpoints and 10,000 paired biological-unit bootstrap
resamples. Candidate sources are CATH 4.4 \citep{waman2025cath44}, PRING
\citep{zheng2025pring}, FLIP2 \citep{didi2026flip2}, MegaScale
\citep{tsuboyama2023megascale}, CAID3 \citep{mehdiabadi2026caid3}, and CAFA5
\citep{depaolis2026cafa5}. It is a qualification profile, not a new aggregate.

\begin{table}[H]
  \centering
  \scriptsize
  \setlength{\tabcolsep}{3pt}
  \begin{tabularx}{\textwidth}{@{}YrrrY@{}}
    \toprule
    Candidate / metric & ESMC-300M & ESMC-600M & ESMC-6B & Decision \\
    \midrule
    CATH 4.4 macro-H mAP & 0.1170 & 0.1067 & 0.1153 & diagnostic; not scale-monotone \\
    PRING Human AUPRC & 0.7035 & 0.7001 & 0.7113 & provisional; intervals failed centering \\
    FLIP2 hierarchical score & 0.7823 & 0.7581 & 0.7805 & diagnostic; limited landscapes \\
    MegaScale median-parent Spearman & 0.6835 & 0.7267 & 0.7207 & qualification candidate \\
    CAID3 macro-protein AP & 0.7292 & 0.7341 & 0.7340 & temporal diagnostic \\
    CAFA5 molecular function & blocked & blocked & blocked & 110 hard proteins $<500$ gate \\
    \bottomrule
  \end{tabularx}
  \caption{P-CORE Next alpha results. Missing tasks are not imputed and no
  aggregate or released-model ranking is issued.}
  \label{tab:pcore-next}
\end{table}

No runnable candidate is strictly monotone across all three scales. MegaScale
is the only new task with a clear paired 6B-over-300M difference:
$+0.0372$ with a 95\% interval of [0.0064, 0.1171]. PRING has only one
negative realization and non-centered intervals. CAFA5 failed its frozen
minimum-population gate. CATH, FLIP2 shift, and CAID3 remain useful
diagnostics but do not yet qualify as promotion gates.

\section{Limitations and next work}

The public corpus reproduces ESMC's source roles and processing logic, but its
OMG/IMG surrogate is much smaller than the reported JGI cluster pool and
cannot reproduce ESMC's unavailable cluster-member sampling. The released
contact protocol matches the ESMC description, snapshot date, and population
size, but the unpublished Biohub chain manifest prevents an item-for-item
identity claim. The current one-hour AutoResearch improvement is single-run
evidence. The historical duration runs use a superseded corpus. The full
released ESMC-6B contact value is still pending. P-CORE Next has not qualified
a successor aggregate.

The immediate work is therefore:

\begin{enumerate}
  \item complete the full 6B contact measurement;
  \item repeat the one-hour incumbent and strongest candidate across seeds;
  \item requalify the four current P-CORE tasks under the same gates applied
        to new candidates;
  \item complete MegaScale confound, weak-model, and probe-stability controls;
  \item repair PRING with a preregistered uncertainty estimator and ten
        negative realizations; and
  \item shadow any P-CORE successor beside the current four-task panel before
        changing model-selection policy.
\end{enumerate}

Nano-ESMC should be read as a reproducible research system and a measured
hill-climbing baseline. It does not yet claim an exact ESMC data reproduction,
a reproduction-quality checkpoint, or a replacement P-CORE suite.

\begingroup
\footnotesize
\setlength{\bibsep}{1pt}
\bibliographystyle{plainnat}
\bibliography{references}
\endgroup

\end{document}
